\documentclass[11pt,twoside]{article}

\oddsidemargin=0in
\evensidemargin=0in
\textwidth=6.5in
\topmargin=0in
\textheight=8.5in
\parindent=0in
\parskip=0in

\usepackage{amsfonts}
\usepackage{tabularx}
\usepackage{amsmath}
\usepackage{amssymb}
\usepackage{graphicx}
\usepackage{mathtools}
%\usepackage{boxedminipage}
\usepackage{colortbl}
\usepackage{comment}

\usepackage{epsfig}
\usepackage{fancyhdr}
\usepackage{layout}
\pagestyle{fancy}
%\headheight = 14pt
\usepackage{hyperref}
\usepackage{multicol}
\usepackage{soul}
\usepackage{pgfplots}
\pgfplotsset{compat=1.16}
\usepackage{tikz}
\usetikzlibrary{external}
\tikzexternalize[prefix=figures/]
\usepackage[inline]{enumitem}
\usepackage{filecontents}

\usetikzlibrary{shapes.geometric}
\usetikzlibrary{arrows}
\usetikzlibrary{patterns}

% Circled numbers
\newcommand*\circled[1]{\tikz[baseline=(char.base)]{
    \node[shape=circle,draw,inner sep=2pt] (char) {#1};}}
    
\DeclareRobustCommand{\inlinegraphic}[1]{%
  \begingroup\normalfont
  {\huge\includegraphics[height=\fontcharht\font`\B]{#1}}%
  \endgroup
}

% Augmented matrices
\makeatletter
\renewcommand*\env@matrix[1][*\c@MaxMatrixCols c]{%
  \hskip -\arraycolsep
  \let\@ifnextchar\new@ifnextchar
  \array{#1}}
\makeatother



\newcommand{\p}{{\mathbb P}}
\newcommand{\dydx}{\frac{dy}{dx}}
\newcommand{\blank}[1]{\underline{\hspace{#1}}}
\newcommand{\dsst}{\displaystyle}
\newcommand{\dd}[2]{\displaystyle \frac{d{#1}}{d{#2}}}
\newcommand{\makearrow}[2]{\overset{#1}{\underset{#2}{\xrightarrow{\hspace{2cm}}}}}    
\pgfplotsset{compat=1.16}


\newcommand{\bpm}{\begin{pmatrix}}
\newcommand{\epm}{\end{pmatrix}}
\newcommand{\mn}{\llap{-}}

\begin{document}
%\section*{Page 42}
%\begin{enumerate}
%\item[24.] $\bpm[cc|c] 2 & 3 & 1 \\ 4 & 1 & 17\epm \makearrow{r_2-2r_1 \rightarrow r_2}{}\bpm[cc|c] 2 & 3 & 1 \\ 0 & \mn5 & 15\epm$. So the triangular system is $\bpm[cc] 2 & 3 \\ 0 & -5\epm\bpm[c] x_1 \\ x_2 \epm = \bpm[c]1\\15\epm$. $-5x_2 = 15$, so $x_2 = -3$. We get $2x_1  + 3(-3) = 1$, so $x_1=5$.\\
%\item[25.] $\bpm[ccc|c] 1 & 2 & 3 & 1 \\ 2 & 3 & 4 & 2 \\ 3 & 5 & 7 & 6\epm \makearrow{r_2-2r_1 \rightarrow r_2}{r_3-3r_1 \rightarrow r_3} \bpm[ccc|c] 1 & 2 & 3 & 1 \\ 0 & \mn1 & \mn2 & 0 \\ 0 & \mn1 & \mn2 & 3\epm \makearrow{r_3-r_2 \rightarrow r_3}{} \bpm[ccc|c] 1 & 2 & 3 & 1 \\ 0 & \mn1 & \mn2 & 0 \\ 0 & 0 & 0 & 3\epm$. We get $0=3$, so there's no solution. If we change that last 6 to a 3, there is a solution (actually an infinite number of them).\\
%\end{enumerate}

\section*{Page 66}
\begin{enumerate}
  \item[1.] \begin{enumerate}
    \item $\bpm 1 & 0 & 0 \\ \mn5 & 1 & 0 \\ 0 & 0 & 1\epm$
    \item $\bpm 1 & 0 & 0 \\ 0 & 1 & 0 \\ 0 & 7 & 1\epm$
    \item $\bpm 0 & 1 & 0 \\ 0 & 0 & 1 \\ 1 & 0 & 0\epm$
  \end{enumerate} 
  
  \item[3.] $E_{21}=\bpm 1 & 0 & 0 \\ \mn4 & 1 & 0 \\ 0 & 0 & 1\epm$, $E_{31}=\bpm 1 & 0 & 0 \\ 0 & 1 & 0 \\ 2 & 0 & 1\epm$, $E_{32}=\bpm 1 & 0 & 0 \\ 0 & 1 & 0 \\ 0 & \mn2 & 1\epm$, $M=\bpm 1 & 0 & 0 \\ \mn4 & 1 & 0 \\ 10 & \mn2 & 1\epm$.
  
  \item[6.] Let $A=\bpm 1 & 2 & 4 \\ 1 & 2 & 4 \\ 1 & 2 & 4 \epm$, and let $\vec{x}=[1,2,3]^T$. Then $A\vec{x}=[16,16,16]^T$. Elimination produces two rows of zeros, so there is only one pivot.
  
  \item[7.] \begin{enumerate}
    \item To invert that step you should \underline{add} 7 times row \underline{~~1~~} to row \underline{~~3~~}.
    \item $E^{-1}=\bpm 1 & 0 & 0 \\ 0 & 1 & 0 \\ 7 & 0 & 1\epm$
    \item $EE^{-1}=\bpm 1 & 0 & 0 \\ 0 & 1 & 0 \\ -7 & 0 & 1\epm\bpm 1 & 0 & 0 \\ 0 & 1 & 0 \\ 7 & 0 & 1\epm = \bpm 1 & 0 & 0 \\ 0 & 1 & 0 \\ 0 & 0 & 1\epm$
  \end{enumerate}
  
  \item[14.] $E_{21} = \bpm 1 & 0 & 0 & 0 \\ \frac12 & 1 & 0 & 0 \\ 0 & 0 & 1 & 0 \\ 0 & 0 & 0 & 1\epm$, $E_{32} = \bpm 1 & 0 & 0 & 0 \\ 0 & 1 & 0 & 0 \\ 0 & \frac23 & 1 & 0 \\ 0 & 0 & 0 & 1\epm$, $E_{43} = \bpm 1 & 0 & 0 & 0 \\ 0 & 1 & 0 & 0 \\ 0 & 0 & 1 & 0 \\ 0 & 0 & \frac34 & 1\epm$
  
  \item[17.] The equation is $\bpm 1 & 1 & 1 \\ 1 & 2 & 4 \\ 1 & 3 & 9\epm \bpm a \\ b \\ c\epm = \bpm 4 \\ 8 \\ 14\epm$. Augmenting and row reducing, we get $\bpm[ccc|c] 1&1&1&4\\ 0 & 1 & 3 & 4 \\ 0 & 0 & 2 & 2\epm$. By back sub, this gives use $c=1$, $b=1$, $a=2$.
  
  \item[24.] $[A|\vec{u}] = \bpm[cc|c] 2 & 3 & 1 \\ 4 & 1 & 17\epm$ becomes $[U|\vec{c}]=\bpm[cc|c] 2 & 3 & 1 \\ 0 & \mn 5 & 15\epm$. So $\vec{x}=[5,-3]^T$.
  
  \item[25.] $[A|\vec{b}] = \bpm[ccc|c] 1 & 2 & 3 & 1 \\ 2 & 3 & 4 & 2 \\ 3 & 5 & 7 & 6\epm$ reduces to $\bpm[ccc|c] 1 & 2 & 3 & 1 \\ 0 & \mn1 & \mn2 & 0 \\ 0 & 0 & 0 & 3\epm$. This gives $0=3$ as the last equation, so no solution. Changing the 6 to a 3 gives a solution (actually, an infinite number of them -- see later).
  
\end{enumerate}

\section*{Page 82}
\begin{enumerate}
  \item[32.] $AX=I$. $x_1$, $x_2$, and $x_3$ are exactly what you get on the right by augmenting the identity to $A$ and doing Gauss-Jordan.
\end{enumerate}

\section*{Page 92}
\begin{enumerate}
  \item[7.] \begin{enumerate}
    \item eqn 1 + eqn 2 must have right hand side 0, but since row 1 + row 2 = row 3, we get two equations with the same coefficients, but different outputs. Hence no solution.
    \item There is a solution exactly if $b_1+b_2=b_3$.
    \item It becomes $0=b_3$.
  \end{enumerate}
  
  \item[22.] $(A|I) = \bpm[cc|cc]1 & 3 & 1 & 0 \\ 2 & 7 & 0 & 1\epm \makearrow{r_2 - 2r_1 \rightarrow r_2}{}\bpm[cc|cc]1 & 3 & 1 & 0 \\ 0 & 1 & \mn2 & 1\epm\makearrow{r_1-3r_2\rightarrow r_1}{} \bpm[cc|cc]1 & 0 & 7 & \mn3 \\ 0 & 1 & \mn2 & 1\epm = (I|A^{-1})$\\

  \item[26.] $\bpm[cc]1& 2\\2&6\epm\makearrow{r_2-2r_1\rightarrow r_2}{}\bpm[cc]1&2\\0&2\epm \makearrow{\frac{r_2}{2}\rightarrow r_2}{} \bpm[cc]1&2\\0&1\epm \makearrow{r_1-2r_2\rightarrow r_1}{}\bpm[cc]1&0\\0&1\epm$. So $E_{21} = \bpm[cc]1&0\\\mn2&1\epm$, $D^{-1} = \bpm[cc]1&0\\0&\frac12\epm$, and $E_{12} = \bpm[cc]1&\mn2\\0&2\epm$. So $A^{-1} = \bpm[cc]1&\mn2\\0&2\epm\bpm[cc]1&0\\0&\frac12\epm\bpm[cc]1&0\\\mn2&1\epm = \bpm 3 &\mn1 \\ \mn1 & \frac12\epm$.\\

  \item[27.] $\bpm[ccc|ccc] 1 & 0 & 0 & 1 & 0 & 0 \\ 2 & 1 & 3 & 0 & 1 & 0 \\ 0 & 0 & 1 & 0 & 0 & 1\epm \makearrow{r_2 - 2r_1 \rightarrow r_2}{} \bpm[ccc|ccc] 1 & 0 & 0 & 1 & 0 & 0 \\ 0 & 1 & 3 & \mn2   & 1 & 0 \\ 0 & 0 & 1 & 0 & 0 & 1\epm \makearrow{r_2 - 3r_3 \rightarrow r_2}{} \bpm[ccc|ccc] 1 & 0 & 0 & 1 & 0 & 0 \\ 0 & 1 & 0 & \mn2   & 1 & \mn3 \\ 0 & 0 & 1 & 0 & 0 & 1\epm$\\
  \begin{align*}
  \bpm[ccc|ccc] 1  & 1 & 1 & 1 & 0 & 0 \\ 1 & 2 & 2 & 0 & 1 & 0\\ 1 & 2 & 3 & 0 & 0 & 1\epm &\makearrow{r_2 - r_1 \rightarrow r_2}{r_3 - r_1 \rightarrow r_3} \bpm[ccc|ccc] 1 & 1 & 1 & 1 & 0 & 0 \\ 0 & 1 & 1 & \mn1 & 1 & 0 \\ 0 & 1 & 2 & \mn1 & 0 & 1\epm &\makearrow{r_3-r_2 \rightarrow r_3}{}\bpm[ccc|ccc] 1 & 1 & 1 & 1 & 0 & 0 \\ 0 & 1 & 1 & \mn1 & 1 & 0 \\ 0 & 0 & 1 & 0 & \mn 1 & 1\epm\\
   &\makearrow{r_2 - r_3 \rightarrow r_2}{r_1 - r_3 \rightarrow r_2}\bpm[ccc|ccc] 1 & 1 & 0 & 1 & 1 & \mn 1\\ 0 & 1 & 0 & \mn1 & 2 & \mn 1 \\ 0 & 0 & 1 & 0 & \mn1 & 1\epm &\makearrow{r_1 - r_2 \rightarrow r_1}{} \bpm[ccc|ccc] 1 & 0 & 0 & 2 & \mn1 & 0\\ 0 & 1 & 0 & \mn1 & 2 & \mn 1 \\ 0 & 0 & 1 & 0 & \mn1 & 1\epm 
  \end{align*}
\item[30.] $A = \bpm[ccc] a & b & b \\ a & a & b \\ a & a & a\epm \makearrow{r_2-r_1 \rightarrow r_2}{r_3-r_1 \rightarrow r_3}\bpm[ccc] a & b & b \\ 0 & a-b & 0 \\ 0 & a-b & a-b\epm \makearrow{r_3 - r_2 \rightarrow r_3}{} \bpm[ccc] a & b & b \\ 0 & a-b & 0 \\ 0 & 0 & a-b\epm$. So if $a\neq 0$ and $a\neq b$, we have non-zero pivots $a$, $a-b$, and $a-b$. So the matrix is invertible.\\

If $c=0$, we have a row of all zeros, so $C$ is not invertible. If $c=2$, the first two rows are the same, so $C$ is not invertible. If $c=7$, the second and third columns are the same, so $C$ is also not invertible (at this stage, students might have to argue that if $c=7$, the transpose of $A$ isn't invertible, so neither is $A$).
\end{enumerate}

\section*{Page 102}
\begin{enumerate}
\item[5.] $A = \bpm 2 & 1 & 0 \\ 0 & 4 & 2 \\ 6 & 3 & 5\epm \makearrow{r_3 - 3r_1 \rightarrow r_3}{} \bpm 2 & 1 & 0 \\ 0 & 4 & 2 \\ 0 & 0 & 5\epm = U$. $E = \bpm 1 & 0 & 0 \\ 0 & 1 & 0 \\ \mn3 & 0 & 1\epm$. So $L = \bpm 1 & 0 & 0 \\ 0 & 1 & 0 \\ 3 & 0 & 1\epm$.\\

\item[6.] $A = \bpm 1 & 1 & 1 \\ 2 & 4 & 5 \\ 0 & 4 & 0 \epm \makearrow{r_2-2r_1\rightarrow r_2}{}\bpm 1 & 1 & 1 \\ 0 & 2 & 3 \\ 0 & 4 & 0 \epm \makearrow{r_3 -2r_2 \rightarrow r_3}{} \bpm 1 & 1 & 1 \\ 0 & 2 & 3 \\ 0 & 0 & \mn6 \epm = U$. $E_{21} = \bpm 1&0&0\\ \mn2&1&0 \\ 0&0&1\epm$, and $E_{32} = \bpm 1&0&0 \\ 0&1&0 \\ 0&\mn2&1\epm$. So $L=\bpm 1&0&0 \\ 2&1&0 \\ 0&2&1\epm$.\\

\item[7.] $A = \bpm 1 & 0 & 1 \\ 2 & 2 & 2 \\ 3 & 4 & 5 \epm \makearrow{r2-2r_1 \rightarrow r_2}{r3-3r_1 \rightarrow r_3} \bpm 1 & 0 & 1 \\ 0 & 2 & 0 \\ 0 & 4 & 2 \epm
\makearrow{r_3 - 2r_2\rightarrow r_3}{} \bpm 1 & 0 & 1 \\ 0 & 2 & 0 \\ 0 & 0 & 2 \epm = U$. So $E_{21} = \bpm 1 & 0 & 0 \\ \mn2 & 1 & 0 \\ 0 & 0 & 1\epm$, $E_{31} = \bpm 1 & 0 & 0 \\ 0 & 1 & 0 \\ \mn3 & 0 & 1\epm$, $E_{32} = \bpm 1 & 0 & 0 \\ 0 & 1 & 0 \\ 0 & \mn2 & 1\epm$. So $L = \bpm 1 & 0 & 0 \\ 2 & 1 & 0 \\ 3 & 2 & 1\epm$\\

\item[8.]\begin{enumerate}
\item  $E_{32}E_{31}E_{21} = \bpm 1&0&0\\ 0&1&0 \\ 0&\mn c&1\epm\bpm 1&0&0\\ 0&1&0 \\ \mn b&0&1\epm \bpm 1&0&0\\ \mn a&1&0 \\ 0&0&1\epm = \bpm 1&0&0\\ 0&1&0 \\ 0&\mn c&1\epm\bpm 1 & 0 & 0 \\ \mn a & 1 & 0 \\ \mn b & 0 & 1\epm = \bpm 1 & 0 & 0 \\ \mn a & 1 & 0 \\ ac-b & \mn c& 1\epm  = E$\\ 
\item $E_{21}^{-1}E_{31}^{-1}E_{32}^{-1} = \bpm 1&0&0\\ a&1&0 \\ 0&0&1\epm\bpm 1&0&0\\ 0&1&0 \\ b&0&1\epm \bpm 1&0&0\\  0&1&0 \\ 0&c&1\epm = \bpm 1&0&0\\ a&1&0 \\ 0&0&1\epm\bpm 1 & 0 & 0 \\ 0 & 1 & 0 \\ b & c & 1\epm  = \bpm 1&0&0\\ a&1&0 \\ b&c&1\epm = L = A$.
\end{enumerate}

\item[15.] $\bpm 1 & 0 \\ 4 & 1 \epm \bpm c_1 \\ c_2 \epm = \bpm 2 \\ 11\epm$. We get $c_1 = 2$ and $4c_1 + c_2 = 11$, so $8+c_2=11$, or $c_2=3$. Then $\bpm 2 & 4 \\ 0 & 1\epm \bpm x_1 \\ x_2 \epm = \bpm 2 \\ 3\epm$. We get $x_2 = 3$ and $2x_1 + 4x_2 = 2$, or $2x_1 + 12 = 2$, so $x_1 = -5$.\\

$\bpm 1 & 0 \\ 4 & 1 \epm \bpm 2 & 4 \\ 0 & 1\epm = \bpm 2 & 4 \\ 8 & 17 \epm$.\\

So consider the augmented matrix $\bpm[cc|c]2 & 4 & 2\\ 8 & 17 & 11 \epm \makearrow{r_2 - 4r_1 \rightarrow r_2}{} \bpm[cc|c] 2 & 4 & 2 \\ 0 & 1 & 3\epm$. This gives $x_2 =3$ and $2x_1+4x_2 = 2$, or $2x_1 + 12 = 2$, so $x_1 = -5$, as before.\\

\pagebreak

\item[16.] $\bpm 1&0&0 \\ 1&1&0 \\ 1&1&1\epm \bpm c_1\\ c_2\\ c_3\epm = \bpm 4 \\ 5 \\ 6 \epm$. So $c_1 =4$. We get $c_1 + c_2 = 5$, or $4 + c_2 = 5$, so $c_2 = 1$; and $c_1+c_2+c_3 = 6$, so $4 + 1 + c_3 = 6$, or $c_3=1$.\\

$\bpm 1  &  1 & 1\\ 0&1&1 \\ 0&0&1 \epm\bpm x_1 \\ x_2 \\ x_3 \epm =  \bpm 4\\1\\1\epm$. So $x_3 = 1$. We get $x_2+x_3 = 1$, or $x_2 + 1 = 1$, so $x_2 = 0$; and $x_1+x_2+x_3 = 4$, or $x_1 + 0 + 1 = 4$, so $x_1 =3$. \\

$A=\bpm 1&0&0 \\ 1&1&0 \\ 1&1&1\epm \bpm 1  &  1 & 1\\ 0&1&1 \\ 0&0&1 \epm = \bpm 1 & 1 & 1 \\ 1 & 2 & 2 \\ 1 & 2 & 3\epm$
\end{enumerate}

%\section*{Extra Problems}
%\begin{enumerate}
%\item \begin{enumerate}
%        \item \begin{align*} 
%                 2b + c &= -1 \\
%                2a +b + c &= -8 \\
%                a +3b + c  &= 0
%                \end{align*}
%        \item $\bpm 0 & 2 & 1 \\ 2 & 1 & 1 \\ 1 & 3 & 1\epm \bpm a \\ b \\ c\epm = \bpm[c] \mn 1\\ \mn 8 \\ 0\epm$\\
%        \item \begin{align*} \bpm[ccc|c] 0 & 2 & 1 & \mn 1 \\ 2 & 1 & 1 & \mn 8 \\ 1 & 3 & 1 & 0 \epm  &\makearrow{r_1 \longleftrightarrow r_2}{} \bpm[ccc|c] 2 & 1 & 1 & \mn 8 \\ 0 & 2 & 1 & \mn 1 \\ 1 & 3 & 1 & 0 \epm &\makearrow{r_3 - \frac12 r_1 \rightarrow r_3}{} \bpm[ccc|c] 2 & 1 & 1 & \mn 8 \\ 0 & 2 & 1 & \mn 1 \\ 0 & \frac52 & \frac12 & 4 \epm\\
%        &\makearrow{r_3 - \frac54 r_2 \rightarrow r_3}{} \bpm[ccc|c] 2 & 1 & 1 & \mn 8 \\ 0 & 2 & 1 & \mn 1 \\ 0 & 0 & \mn\frac34 & \frac{21}{4} \epm
%        \end{align*}
        
 %       So $c = \frac{\frac{21}4}{-\frac34} = -7$. We get $2b+c = -1$, or $2b-7=-1$, so $b=3$; and $2a+b+c=-8$, or $2a+3-7=-8$, so $a=-2$.
  %      \item The plane is $-2x+3y-7 = z$. 
 %     \end{enumerate}

%\pagebreak

%\item $A=\begin{pmatrix}[ccc] \llap{-}3 & 13 & 38 \\ 3 & \llap{-}2  & 4  \\ 2 & \llap{-}7 & \llap{-}19 \end{pmatrix} \makearrow{r_2 + r_1 \rightarrow r_2}{r_3 + \frac23 r_1}{} %\bpm[ccc] \mn3 & 13 & 38 \\ 0 & 11 & 42 \\ 0 & \frac53 & \frac{19}{3}\epm \makearrow{r_3 - \frac{5}{33} r_2 \rightarrow r_3}{} \bpm[ccc] \mn3 & 13 & 38 \\ 0 & 11 & 42 \\ 0 & 0 & \mn\frac{1}{33}\epm = U$\\

%So $E_1 = \bpm[ccc] 1 & 0 & 0 \\ 1 & 1 & 0 \\ 0&0&1\epm$, $E_2 = \bpm[ccc]1&0&0 \\ 0 & 1 & 0 \\ \frac23 & 0 & 1\epm$, and $E_3 = \bpm 1&0&0 \\ 0&1&0 \\ 0&\mn\frac5{33} & 1\epm$.\\

%We also get $L = \bpm[ccc] 1 & 0 & 0 \\ \mn1 & 1 & 0 \\ \mn\frac23 & \frac5{33} & 1\epm$.
%\end{enumerate}

\end{document}
